\documentclass[aps,prb,onecolumn,nofootinbib,superscriptaddress]{revtex4-2}
\usepackage{amsmath,amssymb,amsthm,mathtools,bm}
\usepackage{booktabs}
\usepackage[colorlinks=true,linkcolor=blue,citecolor=blue,urlcolor=blue]{hyperref}
\usepackage{microtype}
\newcommand{\id}{\mathbf{1}}
\newcommand{\tr}{\operatorname{tr}}
\newcommand{\Tr}{\operatorname{Tr}}
\newcommand{\bk}{\boldsymbol k}
\newcommand{\br}{\boldsymbol r}
\newcommand{\C}{\mathcal C}
\newcommand{\BZ}{\mathrm{BZ}}
\newcommand{\rec}{\boldsymbol{\mathrm m}}
\begin{document}
\title{Band mixing, finite-range Wannier measurements, and the topology of conditioned states}
\author{Class-A fermionic adaptive-circuit project}
\noaffiliation
\date{October 5, 2026}
\begin{abstract}
We distinguish the target occupied-band projector, the frame operator of truncated overcomplete Wannier modes, and the occupied projector of a conditioned Gaussian state. Smooth coherent band mixing preserves the Chern number when the deformed occupied line has nonzero overlap with the target line everywhere in momentum space. We give an explicit interpolation and apply a related norm bound to the auxiliary Hamiltonian generated by finite-range Wannier measurements. For the square truncation used in the manuscript, transforming the bra and ket separately gives a quadratic frame operator; the upper- and lower-band contributions then simplify to a linear function of the truncated Fourier series of the target projector. We record a numerical overlap check at range one and explain why this band result does not establish convergence of the stochastic circuit. We also discuss the compact-support obstruction of Dubail and Read and the distinct issue of particle-number changes.
\end{abstract}
\maketitle
\setcounter{tocdepth}{2}
\tableofcontents

\section{Which object carries the topology?}

The adaptive construction measures overcomplete Wannier (OW) modes obtained from the bands of a Chern insulator, truncates their spatial support, and uses feedback to correct their occupations~\cite{BPJ}. Three objects enter this construction. The target band projector specifies the ideal state. The truncated OW functions specify the operations. The correlation matrix of each conditioned trajectory specifies the state actually prepared. These objects coincide only in special limits.

We use the manuscript convention
\begin{equation}
 h(\bk)=\boldsymbol n(\bk)\cdot\boldsymbol\sigma,\qquad
 \boldsymbol n(\bk)=(\sin k_x,\sin k_y,\alpha-\cos k_x-\cos k_y),\qquad
 P_\pm(\bk)=\frac12\left(\id\pm\widehat{\boldsymbol n}(\bk)\cdot\boldsymbol\sigma\right).
 \label{eq:model}
\end{equation}
For a smooth periodic occupied projector $P(\bk)$, define
\begin{equation}
 \C[P]=\frac{1}{2\pi i}\int_{\BZ}d^2k\,
 \tr\left(P[\partial_{k_x}P,\partial_{k_y}P]\right).
 \label{eq:chern}
\end{equation}
This convention gives $\C[P_-]=1$ for $0<\alpha<2$. Quantization follows from the topology of the occupied bundle, and invariance under a smooth fixed-rank deformation follows because an integer cannot change continuously~\cite{Avron1983}. No globally smooth occupied eigenvector is assumed; local eigenvectors on overlapping momentum patches suffice.

For a pure number-conserving Gaussian trajectory, $G_{ij}=\Tr(\hat\rho\,\hat c_i^\dagger\hat c_j)$ and its real-space occupied projector is $G^{\mathsf T}$. We retain $G_c=2G-\id$. A typical trajectory is not translation invariant, so it need not admit a separate rank-one occupied projector at each momentum. Averaging over records generally destroys the projector property: the topology of a conditioned state and a quantity evaluated on $\overline G$ must not be interchanged. Likewise, a finite-disk Chern estimator or a local marker need not equal an integer at finite size, even when it diagnoses a quantized bulk phase~\cite{BiancoResta2011}.

\section{Coherent band mixing and an explicit interpolation}

Let $P(\bk)$ and $R(\bk)$ be smooth periodic rank-one orthogonal projectors. Here $R$ is a comparison projector, not a reference-system label. On a local momentum patch, write its occupied state as
\begin{equation}
 |v\rangle=a|u_-\rangle+b|u_+\rangle,\qquad
 |a|^2+|b|^2=1,\qquad
 P=|u_-\rangle\langle u_-|,\qquad R=|v\rangle\langle v|.
\end{equation}
The leakage relative to the target band is
\begin{equation}
 \eta(\bk)=\tr[(\id-P)R]=|b(\bk)|^2,\qquad
 \tr(PR)=1-\eta=|a(\bk)|^2.
 \label{eq:overlap}
\end{equation}
The decomposition into two bands alone imposes no constraint: $a$ could vanish. Suppose instead that $\eta(\bk)<1$ at every momentum. Compactness of the Brillouin torus and continuity then give a positive minimum of $1-\eta$.

To remove the upper-band component continuously, define
\begin{equation}
 B_s=P+(1-s)(\id-P),\qquad
 R_s=\frac{B_sRB_s}{\tr(B_sRB_s)},\qquad 0\le s\le1.
 \label{eq:overlap-path}
\end{equation}
Since $R$ has rank one, the numerator is a rank-one positive operator and division by its trace makes it an orthogonal projector. Its denominator is
\begin{equation}
 \tr(B_sRB_s)=1-\eta+(1-s)^2\eta>0.
\end{equation}
Thus $R_s$ is smooth and periodic throughout the interpolation. At $s=0$ it equals $R$, while $PRP=\tr(PR)P$ gives $R_1=P$. It follows that
\begin{equation}
 \boxed{\tr[P(\bk)R(\bk)]>0\ \text{for all }\bk
 \quad\Longrightarrow\quad \C[R]=\C[P].}
 \label{eq:overlap-result}
\end{equation}
Equivalently, projection onto the target band supplies an everywhere invertible map between the occupied line bundles. Different Chern numbers therefore require an orthogonality point. The converse does not hold: an orthogonality point does not by itself imply different topology.

This result allows substantial coherent mixing. It is a condition on the geometry of the occupied subspace, not on closeness of all observables or many-body fidelity. In contrast, a small momentum-averaged value of $\eta$ does not exclude an orthogonality point or a texture concentrated in a small momentum region. For several occupied bands, the corresponding overlap map must have full rank everywhere; a positive trace alone is insufficient.

Caio, Cooper, and Bhaseen provide a physical example in the quenched Haldane model: both post-quench bands become occupied, while the pure-state Chern number is preserved by smooth momentum-dependent unitary evolution~\cite{Caio2015}. Their result concerns a coherent state. It neither equates that state with the final Hamiltonian's ground state nor guarantees quantization of its nonequilibrium Hall response~\cite{Caio2016}. It also does not prove stability under a sequence of conditioned measurements.

\section{Fourier transforming the operator and the OW modes}

\subsection{Both operator indices must be transformed}

Suppressing orbital indices, let $\langle\br|\bk\rangle=N_{\mathrm{uc}}^{-1/2}e^{i\bk\cdot\br}$. An operator transforms as
\begin{equation}
 P(\bk,\bk')=\frac{1}{N_{\mathrm{uc}}}
 \sum_{\br,\br'}e^{-i\bk\cdot\br}P_{\br,\br'}e^{i\bk'\cdot\br'}.
 \label{eq:double-fourier}
\end{equation}
For a translation-invariant kernel $P_{\br,\br'}=P(\br-\br')$, the center-coordinate sum yields $\delta_{\bk,\bk'}$, leaving
\begin{equation}
 P(\bk,\bk')=\delta_{\bk,\bk'}P(\bk),\qquad
 P(\bk)=\sum_{\boldsymbol d}P(\boldsymbol d)e^{-i\bk\cdot\boldsymbol d}.
 \label{eq:relative-fourier}
\end{equation}
The single Fourier series is therefore a consequence of translation invariance, not a transformation of only one operator index.

For a single localized mode, define its unnormalized Fourier amplitude by $W(\bk)=\sum_{\br}e^{-i\bk\cdot\br}W(\br)$. Transforming its rank-one projector instead gives
\begin{equation}
 \langle\bk|W\rangle\langle W|\bk'\rangle
 =\frac{1}{N_{\mathrm{uc}}}W(\bk)W^\dagger(\bk').
 \label{eq:mode-fourier}
\end{equation}
This is not momentum diagonal. Summing the projectors of all lattice translates of $W$ makes it diagonal and gives the Bloch matrix $W(\bk)W^\dagger(\bk)$. This distinction is essential when constructing the truncated frame operator.

\subsection{Expanding the ket and bra into real-space coefficients}

The same construction can be read directly from the momentum-dependent spinor. Expand the ket and bra separately,
\begin{equation}
 |v(\bk)\rangle=\sum_{\br}e^{-i\bk\cdot\br}|w(\br)\rangle,\qquad
 \langle v(\bk)|=\sum_{\br'}e^{i\bk\cdot\br'}\langle w(\br')|.
 \label{eq:spinor-expansion}
\end{equation}
Here $|w(\br)\rangle$ is an orbital-valued Fourier coefficient, not a position basis ket. For a topologically nontrivial band, a normalized eigenvector requires a gauge choice and is not assumed globally smooth. The expansion also applies directly to the globally defined, generally unnormalized OW spinors $P_-(\bk)\tau_\nu$.

The outer product contains two independent real-space sums:
\begin{equation}
 M_v(\bk)\equiv|v(\bk)\rangle\langle v(\bk)|
 =\sum_{\br,\br'}e^{-i\bk\cdot(\br-\br')}
 |w(\br)\rangle\langle w(\br')|.
 \label{eq:outer-double-sum}
\end{equation}
If $v$ is normalized at each momentum, $M_v$ is the occupied projector $R$; otherwise it is simply the outer-product matrix. Introduce the difference coordinate $\boldsymbol d=\br-\br'$. Regrouping Eq.~\eqref{eq:outer-double-sum} gives
\begin{equation}
 M_v(\bk)=\sum_{\boldsymbol d}e^{-i\bk\cdot\boldsymbol d}M_v(\boldsymbol d),\qquad
 M_v(\boldsymbol d)=\sum_{\br'}|w(\br'+\boldsymbol d)\rangle\langle w(\br')|.
 \label{eq:autocorrelation}
\end{equation}
Thus the relative-coordinate kernel is an autocorrelation of the Fourier coefficients. Equivalently, the translation-invariant operator obtained by summing all translates of the real-space mode has kernel
\begin{equation}
 (M_v)_{\br,\br'}=\sum_{\boldsymbol x}
 |w(\br-\boldsymbol x)\rangle\langle w(\br'-\boldsymbol x)|
 =M_v(\br-\br')
 =\int_{\BZ}\frac{d^2k}{(2\pi)^2}
 e^{i\bk\cdot(\br-\br')}M_v(\bk).
 \label{eq:translation-sum}
\end{equation}
This displays explicitly how the separate ket and bra expansions, the double real-space sum, and the single momentum integral are consistent. On a finite periodic lattice the integral is replaced by $N_{\mathrm{uc}}^{-1}\sum_{\bk}$.

Truncating the wavefunction gives $w_w(\br)=g_w(\br)w(\br)$, so its autocorrelation becomes
\begin{equation}
 M_{v,w}(\boldsymbol d)=\sum_{\br'}
 g_w(\br'+\boldsymbol d)g_w^*(\br')
 |w(\br'+\boldsymbol d)\rangle\langle w(\br')|.
 \label{eq:window-autocorrelation}
\end{equation}
Both window factors are required. This is generally not $g_w(\boldsymbol d)M_v(\boldsymbol d)$. A square wavefunction window of half-width $w$ gives an outer-product kernel supported within $|d_x|,|d_y|\le2w$, before periodic wrapping. A real-space normalization of the truncated mode divides this kernel by its squared norm. Normalization separately at every momentum is a different operation and can restore long-range coefficients. In particular, $M_{v,w}$ is not automatically an occupied-band projector and must not be confused with the spectral projector $R_w$ defined below.

\subsection{The window truncates a wavefunction before its outer product}

Let $g_w(\boldsymbol d)$ be real, inversion symmetric, equal to one at the origin, and of finite support. For the manuscript it is the indicator of $|d_x|,|d_y|\le w$. Define
\begin{equation}
 F_w(\bk)=\sum_{\boldsymbol d}g_w(\boldsymbol d)P_-(\boldsymbol d)e^{-i\bk\cdot\boldsymbol d}.
 \label{eq:F}
\end{equation}
Hermiticity of the original kernel and inversion symmetry of the window imply $F_w^\dagger=F_w$. Since $P_+=\id-P_-$ and the window retains the origin, its upper-band counterpart is $\id-F_w$. The normalized truncated OW amplitudes are
\begin{equation}
 W_{\nu,-;w}(\bk)=\frac{F_w(\bk)\tau_\nu}{\sqrt{Z_{\nu,-;w}}},\qquad
 W_{\nu,+;w}(\bk)=\frac{[\id-F_w(\bk)]\tau_\nu}{\sqrt{Z_{\nu,+;w}}},
 \label{eq:OW}
\end{equation}
where $Z_{\nu,-;w}=\sum_{\boldsymbol d}|g_w(\boldsymbol d)|^2\|P_-(\boldsymbol d)\tau_\nu\|^2$ and similarly for the upper band. These are real-space normalization factors, not separate normalizations at every momentum.

Writing $D_{\pm,w}=\sum_\nu\tau_\nu\tau_\nu^\dagger/Z_{\nu,\pm;w}$, the translation-summed frame operators are
\begin{equation}
 S_{-,w}=F_wD_{-,w}F_w^\dagger,\qquad
 S_{+,w}=(\id-F_w)D_{+,w}(\id-F_w)^\dagger.
 \label{eq:frames-general}
\end{equation}
In particular, neither $S_{-,w}$ nor the projector of an individual truncated mode should be identified with $F_w$. Windowing a wavefunction and then forming its outer product differs from windowing an operator kernel.

For the model in Eq.~\eqref{eq:model}, the square window, and $\tau_{A,B}=(1,\pm1)^{\mathsf T}/\sqrt2$, all four $Z_{\nu,\pm;w}$ are equal to a common $Z_w$. To see this, write $F_w=(\id-\boldsymbol d_w\cdot\boldsymbol\sigma)/2$. The Brillouin-zone average of $d_{w,x}$ vanishes by reflection symmetry, and the two trial orbitals have zero $\sigma^y$ and $\sigma^z$ expectation values. Averaging $\tau_\nu^\dagger F_w^2\tau_\nu$, or its upper-band counterpart, therefore gives the same $(1+\langle|\boldsymbol d_w|^2\rangle_{\BZ})/4$.

Completeness of the two trial orbitals then yields
\begin{equation}
 S_{-,w}=\frac{F_w^2}{Z_w},\qquad
 S_{+,w}=\frac{(\id-F_w)^2}{Z_w},\qquad
 h_w\equiv S_{+,w}-S_{-,w}=\frac{\id-2F_w}{Z_w}.
 \label{eq:frame-hamiltonian}
\end{equation}
The last equality follows from cancellation of the quadratic terms. It is not obtained by omitting the bra transform. For other trial orbitals or windows with unequal normalization factors, Eq.~\eqref{eq:frames-general} must be retained.

\section{A sufficient truncation bound and its numerical check}

The negative-energy band of Eq.~\eqref{eq:frame-hamiltonian} is the spectral subspace of $F_w$ above $1/2$. Denote its projector by $R_w$. A sufficient condition for this band to retain the target topology is
\begin{equation}
 \delta_w\equiv\sup_{\bk}\|F_w(\bk)-P_-(\bk)\|_{\mathrm{op}}<\frac12.
 \label{eq:delta}
\end{equation}
Indeed, along $F_s=P_-+s(F_w-P_-)$, the two eigenvalues stay within $s\delta_w$ of zero and one. Neither crosses $1/2$. The spectral projector $R_s$ above $1/2$ is therefore a smooth fixed-rank interpolation from $P_-$ to $R_w$, proving
\begin{equation}
 \boxed{\delta_w<\frac12\quad\Longrightarrow\quad\C[R_w]=\C[P_-].}
 \label{eq:truncation-result}
\end{equation}
The endpoint Hamiltonian also obeys $\min_{\bk,n}|E_n[h_w]|\ge(1-2\delta_w)/Z_w$. Equivalently, $\id-2F_s$ connects it, up to a positive normalization, to the spectrally flattened target Hamiltonian without closing the gap. This conclusion is an application of band homotopy~\cite{Avron1983}; it is not a statement about relaxation of the measurement channel.

For an indicator window, the triangle inequality gives
\begin{equation}
 \delta_w\le\sum_{\boldsymbol d\notin\mathrm{window}}
 \|P_-(\boldsymbol d)\|_{\mathrm{op}}.
 \label{eq:tail}
\end{equation}
At a fixed gapped target parameter, exponential decay of the projector kernel makes this tail vanish for increasing windows. Thus sufficiently large symmetric truncations satisfy Eq.~\eqref{eq:delta}. For a general real symmetric window the bound instead contains $|1-g_w(\boldsymbol d)|$ in every summand. If unequal normalizations prevent Eq.~\eqref{eq:frame-hamiltonian}, one may instead compare the full Hermitian $h_w$ with its untruncated limit $h_\infty$: the sufficient condition $\sup_{\bk}\|h_w-h_\infty\|_{\mathrm{op}}<\min_{\bk,n}|E_n[h_\infty]|$ keeps their straight interpolation gapped at zero. A small discarded squared norm or integrated opposite-band weight is a different, weaker diagnostic and must not be substituted for the uniform operator-norm bound. Near a target gap closing the necessary range can grow.

We checked the manuscript's shortest square window, $\alpha=1,w=1$, by reconstructing the Fourier coefficients on a $1024\times1024$ grid and evaluating the finite Fourier sum on two independent momentum grids. The companion script records the following values; its eigenvalue and overlap calculations do not run the circuit.
\begin{table}[t]
\caption{Finite-grid checks for the auxiliary band. The maximum error and minimum overlap are sampled extrema, not certified continuum bounds. The energy entry is the minimum absolute eigenvalue, rather than twice that value.}
\begin{ruledtabular}
\begin{tabular}{lcc}
 & $201\times201$ & $401\times401$\\
\hline
Maximum $\|F_w-P_-\|_{\mathrm{op}}$ & $0.09873199$ & $0.09875389$\\
Minimum $\tr(P_-R_w)$ & $0.99019725$ & $0.99018986$\\
Minimum $|E[h_w]|$ & $1.80242607$ & $1.80242648$\\
Common $Z_w$ & $0.49607752$ & $0.49607752$
\end{tabular}
\end{ruledtabular}
\label{tab:numerics}
\end{table}
The independently assembled quadratic frame expression agrees with Eq.~\eqref{eq:frame-hamiltonian} to better than $3\times10^{-15}$ in maximum matrix-entry magnitude. The numerical margin to $\delta_w=1/2$ is large, and the sampled lower-band overlap exceeds $0.9901$. However, finite momentum grids and numerical Fourier coefficients do not constitute a rigorous continuum certificate. The theorem is conditional on Eq.~\eqref{eq:delta}; the table is a numerical application of that condition. The auxiliary gap and overlap are not measurements of the occupied projector of a stochastic trajectory.

\subsection{Single-site truncation: where the two proofs fail}
\label{sec:single-site}

The limit $w=0$ retains only the central unit cell, including its two orbitals. It is therefore an onsite orbital measurement, not necessarily measurement of one physical orbital. The truncated matrix is independent of momentum. Reflection symmetry cancels the averages of $\widehat n_x$ and $\widehat n_y$, giving
\begin{equation}
 F_0=\int_{\BZ}\frac{d^2k}{(2\pi)^2}P_-(\bk)
 =\frac12(\id-d_0\sigma^z),\qquad
 d_0=\int_{\BZ}\frac{d^2k}{(2\pi)^2}\widehat n_z(\bk).
 \label{eq:onsite-frame}
\end{equation}
For $\alpha=1$, numerical quadrature gives $d_0\simeq0.49116916$. The two trial orbitals have equal normalization $Z_0=(1+d_0^2)/4$, so the auxiliary Hamiltonian and its negative-energy projector are
\begin{equation}
 h_0=\frac{d_0}{Z_0}\sigma^z,\qquad
 R_0=|\downarrow\rangle\langle\downarrow|,\qquad
 \C[R_0]=0.
 \label{eq:onsite-projector}
\end{equation}
The last equality is exact: a constant projector has zero momentum derivatives. Thus the single-site auxiliary band is gapped but topologically trivial, whereas the target has $\C[P_-]=1$.

The failure of the overlap proof is explicit at $\Gamma=(0,0)$. Here $\boldsymbol n=(0,0,-1)$ and $P_-(\Gamma)=|\uparrow\rangle\langle\uparrow|$, hence
\begin{equation}
 \tr[P_-(\Gamma)R_0]=0,\qquad \eta(\Gamma)=1.
\end{equation}
In Eq.~\eqref{eq:overlap-path}, $B_sR_0B_s=(1-s)^2R_0$ at this momentum. Its normalization vanishes at $s=1$: the path cannot reach the target projector. Merely writing an onsite spinor as a linear combination of the target upper and lower bands does not prevent its lower-band coefficient from vanishing.

The norm argument fails for the same physical reason. At $\Gamma$,
\begin{equation}
 \|F_0-P_-(\Gamma)\|_{\mathrm{op}}
 =\frac{1+d_0}{2}\simeq0.74558458>\frac12.
\end{equation}
More than the sufficient bound is lost: the particular interpolation used in the proof actually closes its gap. With $q=(1+d_0)/2$, its value at $\Gamma$ is
\begin{equation}
 F_s(\Gamma)=(1-s)P_-(\Gamma)+sF_0
 =\begin{pmatrix}1-sq&0\\0&sq\end{pmatrix},\qquad
 s_*=\frac{1}{1+d_0}.
 \label{eq:onsite-crossing}
\end{equation}
At $s=s_*\in(0,1)$ both eigenvalues equal $1/2$, and $\id-2F_s$ vanishes at $\Gamma$. The endpoints can therefore both be gapped while carrying different Chern numbers; no gapped interpolation connects them.

An individual nonzero onsite OW mode likewise has a momentum-independent normalized projector, since its center-dependent Fourier phase cancels between ket and bra. This mode projector need not equal $R_0$, which is defined by the combined frame Hamiltonian. Neither object should be identified with a stochastic trajectory without further dynamical analysis. Independently, operations confined to individual unit cells cannot create intercell entanglement from an initially cell-factorized state. The single-site limit thus removes the spatial structure needed by the preparation protocol, rather than providing a counterexample to either sufficient-condition proof.

\section{Compact support, overcompleteness, and the no-go theorem}

Dubail and Read show that a polynomial and analytic complex filled-band bundle is topologically trivial; for their nontrivial Gaussian tensor-network states, a short-range quadratic parent cannot have a bulk gap. Their Sec.~III G also rules out a set of compactly supported Wannier-type functions spanning an exact nontrivial band at every momentum, even when redundancy is allowed~\cite{DubailRead2015}. Overcompleteness alone therefore does not evade the compact-support obstruction.

The present construction does not supply such a set. The exact OW functions are band confined but have spatial tails. Their truncated versions have compact support but leak into the opposite band and generally cease to span only the intended occupied space. Furthermore, $R_w$ is obtained by spectral projection of $F_w$, a nonlinear operation that generally restores spatial tails. A finite Fourier polynomial $F_w$ can have a topologically nontrivial spectral subspace without supplying polynomial sections that span that subspace. This is why a finite-range, gapped $h_w$ is consistent with the theorem.

For one OW family the lower-band overlap $\langle u_-|W_{\nu,-;w}\rangle$ can vanish. Equation~\eqref{eq:overlap-result} cannot be applied to that mode as though it were the entire occupied band. The two families must first be treated together, as in Eq.~\eqref{eq:frames-general}, and even that frame construction defines an auxiliary band rather than the conditioned state.

\section{Particle number and local gain or loss}

For a Slater determinant with real-space occupied projector $P$, the physical particle number is $N_F=\tr P$. A successful local creation operation $\hat c^\dagger(v)=\sum_i v_i\hat c_i^\dagger$ adds the normalized orbital
\begin{equation}
 |\phi_+\rangle=\frac{(\id-P)|v\rangle}{\sqrt{\langle v|(\id-P)|v\rangle}},\qquad
 P'=P+|\phi_+\rangle\langle\phi_+|.
 \label{eq:gain}
\end{equation}
The occupied component of $v$ is excluded by the Pauli principle. Similarly, a successful annihilation operation removes
\begin{equation}
 |\phi_-\rangle=\frac{P|v\rangle}{\sqrt{\langle v|P|v\rangle}},\qquad
 P'=P-|\phi_-\rangle\langle\phi_-|.
 \label{eq:loss}
\end{equation}
Only nonzero-probability operations are included, so these denominators are positive. Both identities describe pure Slater inputs and the specified gain or loss branch, not an outcome-averaged reset.

These changes alter the rank by one and have $\|P'-P\|_{\mathrm{op}}=1$. They therefore fall outside a small-norm, equal-rank comparison, even if $|\Delta N_F|/N_{\mathrm{uc}}$ is tiny. Nevertheless, a finite number of sufficiently localized added or removed orbitals does not change the thermodynamic bulk Chern index in the usual real-space index framework. The modification is finite rank, and the bulk Fredholm index is stable to the associated compact perturbation when the locality conditions making the index well defined are maintained. Finite-region marker values can still change near the orbitals. The quantum Hall plateau result of Bellissard, van Elst, and Schulz-Baldes gives a stronger setting in which particle density changes while the invariant remains fixed, as the Fermi level moves through a localized spectral region~\cite{Bellissard1994}.

If $v$ is local and $P$ has sufficiently rapid bulk decay, the orbitals in Eqs.~\eqref{eq:gain} and \eqref{eq:loss} inherit localization for each allowed event. This statement does not give bounds uniform over all outcomes or all times. Boundary-localized charge can likewise change without changing a bulk index far from the boundary. Conversely, populating extended bulk states can destroy the insulating assumptions; a partially filled clean band need not define a smooth fixed-rank occupied bundle. A small net particle-number deviation does not distinguish these cases and can conceal many particle--hole excitations.

\section{What is established for the adaptive circuit}

The deterministic results concern the relation between the exact target band and the auxiliary occupied band generated by truncation. They establish neither a common dark state of all truncated resets nor equality of that auxiliary band with the final trajectory projector. The two truncated families can span the full orbital space, so filling all lower-family modes while emptying all upper-family modes is incompatible. Subsequent operations can disturb occupations established earlier. The stationary object is then an ensemble of conditioned states rather than one common stabilized wavefunction.

A proof that every relevant long-time trajectory retains a given bulk index would require control of the trajectory projectors, their spatial decay, and the appropriate limits. Individual records break translation symmetry and may change physical particle number. Repeating local operations for times growing with system size is also not an exact preparation by a fixed-depth circuit or a fixed-bond-dimension tensor network. An averaged-channel relaxation gap alone supplies none of these missing controls.

The present numerical evidence is separate: the manuscript's real-space disk estimator, evaluated in each trajectory before averaging, approaches one in the slab interior. Its cycle-40 mean deviation decreases from approximately $2.4\times10^{-3}$ at $L=20$ to $1.4\times10^{-5}$ at $L=40$. This supports the intended bulk phase on the tested sizes; it is not an exact finite-size quantization theorem or a claim that every trajectory has an identical marker. Likewise, the critical and chiral boundary diagnostics do not identify the conditioned states with a particular equilibrium ground state.

A concise statement suitable for the truncation discussion is: \emph{Finite-range truncation mixes the target bands and frustrates exact OW stabilization. It need not change the topological phase: the auxiliary occupied band remains connected to the target band when the truncation error is uniformly smaller than the spectral separation, and the prepared trajectories independently exhibit a bulk Chern estimator approaching one. The latter observation is a dynamical numerical result, rather than a consequence of small integrated band leakage alone.}

\appendix
\section{Reproduction and source scope}
The accompanying \texttt{check\_truncation.py} reads the existing analytical renderer \path{technical_report/analyze_ow_truncation.py} without calling its campaign or plotting entry point. It computes the deterministic Fourier reconstruction, checks the quadratic frame identity, and writes \texttt{truncation\_check.json}. The receipt records the source hash, grid sizes, NumPy version, sampled extrema, and numerical residuals. The $1024\times1024$ Fourier grid is the existing manuscript analysis convention. No circuit simulation, trajectory resampling, fit, or modification of manuscript assets is involved. The local script path is resolved from the repository root, so the check can be rerun after moving the checkout.

The citations serve different purposes: Ref.~\cite{Avron1983} supplies the band-homotopy background; Refs.~\cite{Caio2015,Caio2016} distinguish coherent state topology from nonequilibrium response; Ref.~\cite{DubailRead2015} gives the compact-support obstruction; Ref.~\cite{Bellissard1994} treats quantization under localization assumptions. None is cited as a proof of convergence of this adaptive trajectory ensemble.
\bibliographystyle{apsrev4-2}
\bibliography{references}
\end{document}
